\chapter{Pendahuluan}
\label{chap:pendahuluan}

\section{Latar Belakang}
\label{sec:latar-belakang}

Penyakit kardiovaskular (\textit{cardiovascular disease}, CVD) merupakan penyebab utama kematian secara global. Menurut data Organisasi Kesehatan Dunia (WHO), diperkirakan 17,9 juta orang meninggal akibat CVD setiap tahunnya, mewakili sekitar 32\% dari seluruh kematian di dunia~\cite{who_cvd}. Angka ini menunjukkan urgensi pengembangan sistem deteksi dini yang dapat membantu mengurangi angka mortalitas akibat penyakit jantung.

Aritmia merupakan salah satu manifestasi CVD yang paling berbahaya karena seringkali terjadi tanpa gejala yang jelas (\textit{asimptomatik}). Tanpa terdeteksi, aritmia dapat berujung pada komplikasi serius seperti \textit{stroke}, gagal jantung, atau kematian jantung mendadak (\textit{sudden cardiac death}). Oleh karena itu, deteksi dini aritmia melalui monitoring elektrokardiogram (ECG) secara kontinu diperlukan agar intervensi medis dapat dilakukan tepat waktu~\cite{sheeba_rani_2026}.

Metode monitoring jantung konvensional, seperti perangkat Holter monitor, memiliki beberapa keterbatasan. Perangkat Holter umumnya mahal, berukuran relatif besar, memerlukan kunjungan ke fasilitas kesehatan untuk pemasangan dan pengambilan data, serta hanya mampu merekam dalam durasi terbatas (24--48 jam). Hal ini menjadikan monitoring jantung kontinu jangka panjang sulit diakses masyarakat luas, terutama di daerah dengan akses kesehatan terbatas.

Perkembangan teknologi Internet of Things (IoT) dan perangkat \textit{wearable} membuka peluang baru untuk monitoring kesehatan yang lebih terjangkau dan aksesibel. Integrasi dengan paradigma \textit{edge computing} memungkinkan pemrosesan data dilakukan langsung di perangkat (\textit{on-device}), mengurangi ketergantungan pada infrastruktur \textit{cloud} dan meminimalkan latensi respons~\cite{tinyml_iot_review, yahyati_2025}.

Namun demikian, mayoritas sistem IoT untuk monitoring kesehatan saat ini masih menganut arsitektur \textit{cloud-centric}, di mana seluruh data sensor dikirim ke server \textit{cloud} untuk diproses. Pendekatan ini berlatensi tinggi untuk aplikasi deteksi aritmia \textit{real-time} dan bergantung pada konektivitas internet yang stabil. Pendekatan yang sama juga membawa potensi masalah privasi data kesehatan serta konsumsi daya dan \textit{bandwidth} yang tinggi. Selain itu, model kecerdasan buatan untuk deteksi aritmia umumnya dievaluasi pada lingkungan PC server tanpa mempertimbangkan batasan yang ada pada mikrokontroler. \textit{Pipeline} data dari \textit{edge} ke \textit{cloud} juga belum dirancang dengan mekanisme resiliensi terhadap gangguan jaringan~\cite{escobedo_2026, yaacoub_2020}.

Berdasarkan permasalahan tersebut, penelitian ini mengajukan solusi berupa implementasi sistem monitoring detak jantung kontinu berbasis IoT \textit{edge computing} menggunakan pendekatan Tiny Machine Learning (TinyML). Sistem ini menjalankan preprocessing sinyal ECG dan inferensi model deteksi aritmia secara langsung pada mikrokontroler ESP32-S3, serta membangun \textit{pipeline} \textit{edge-to-cloud} yang resilien dengan mekanisme bufer lokal menggunakan protokol MQTT. Dengan pendekatan ini, keputusan deteksi dihasilkan di perangkat tanpa menunggu \textit{cloud}, dan data monitoring tetap tersimpan ketika koneksi terputus.

\section{Permasalahan}
\label{sec:permasalahan}

Berdasarkan latar belakang yang telah diuraikan, rumusan masalah dalam penelitian ini diuraikan pada butir berikut.
\begin{enumerate}
  \item Bagaimana merancang algoritma preprocessing sinyal ECG yang ringan dan efisien untuk dijalankan secara \textit{real-time} pada mikrokontroler berdaya rendah?
  \item Bagaimana mengimplementasikan model TinyML teroptimasi untuk deteksi aritmia secara \textit{on-device} pada mikrokontroler ESP32-S3?
  \item Bagaimana membangun \textit{pipeline} IoT \textit{edge-to-cloud} yang resilien dengan mekanisme \textit{buffering} lokal saat terjadi gangguan jaringan?
\end{enumerate}

\section{Tujuan}
\label{sec:tujuan}

Penelitian ini mengajukan pendekatan baru untuk deteksi dini aritmia dengan menjalankan preprocessing sinyal ECG dan inferensi model TinyML secara langsung pada mikrokontroler ESP32-S3, serta membangun \textit{pipeline} \textit{edge-to-cloud} yang resilien sebagai klaim orisinalitas. Secara lebih rinci, tujuan penelitian ini diuraikan pada butir berikut.
\begin{enumerate}
  \item Mengimplementasikan algoritma preprocessing ECG ringan yang mampu berjalan secara \textit{real-time} pada mikrokontroler berdaya rendah.
  \item Menerapkan model TinyML terkompresi (kuantisasi INT8) untuk deteksi aritmia secara \textit{on-device} pada ESP32-S3 tanpa ketergantungan pada komputasi \textit{cloud}.
  \item Membangun \textit{pipeline} \textit{edge-to-cloud} yang resilien menggunakan protokol MQTT dengan mekanisme bufer lokal untuk menjaga kontinuitas data monitoring saat koneksi terputus.
\end{enumerate}

Kriteria keberhasilan setiap tujuan dirumuskan secara terukur pada rencana pengujian di Subbab~\ref{sec:rencana-pengujian}, dan ketercapaiannya dianalisis per tujuan pada Bab~\ref{chap:eksperimen}.

\section{Manfaat}
\label{sec:manfaat}

Manfaat yang diharapkan dari penelitian ini diuraikan pada butir berikut.
\begin{enumerate}
  \item Menghasilkan \emph{prototype} monitoring aritmia \textit{real-time} yang lebih terjangkau dan portabel dibandingkan perangkat medis konvensional.
  \item Membuktikan kelayakan penerapan TinyML pada mikrokontroler ESP32-S3 untuk domain kesehatan, khususnya deteksi aritmia berbasis sinyal ECG.
  \item Menjadi referensi implementasi \textit{edge computing} untuk sistem monitoring biomedis berbasis IoT \textit{open-source} yang dapat dikembangkan lebih lanjut oleh peneliti dan praktisi.
\end{enumerate}

\section{Sistematika Penulisan}
\label{sec:sistematika}

Sistematika penulisan laporan ini disusun sebagai berikut.

\noindent
\textbf{\makebox[1.5cm][l]{BAB 1}PENDAHULUAN}\\
\hspace*{1.5cm} Bab ini berisi latar belakang masalah yang melatarbelakangi penelitian, perumusan masalah, tujuan, manfaat, serta sistematika penulisan laporan secara keseluruhan. Bab ini memberikan gambaran umum mengenai urgensi pengembangan sistem monitoring aritmia berbasis IoT \textit{edge computing} dengan pendekatan TinyML.

\noindent
\textbf{\makebox[1.5cm][l]{BAB 2}KAJIAN PUSTAKA}\\
\hspace*{1.5cm} Bab ini menguraikan teori-teori dasar dan literatur yang relevan dengan topik penelitian, meliputi aritmia dan sinyal ECG, IoT \textit{edge computing}, Tiny Machine Learning (TinyML), 1D-CNN dan kuantisasi model, mikrokontroler ESP32-S3, sensor ECG AD8232, protokol MQTT, platform ThingsBoard, basis data MIT-BIH, serta evaluasi \textit{inter-patient} dan augmentasi data. Selain itu, dibahas pula penelitian terdahulu yang menjadi pembanding untuk menegaskan posisi penelitian ini.

\noindent
\textbf{\makebox[1.5cm][l]{BAB 3}DESKRIPSI SISTEM}\\
\hspace*{1.5cm} Bab ini menjelaskan rancangan sistem secara keseluruhan, termasuk desain arsitektur \textit{edge-to-cloud}, perancangan perangkat keras berbasis ESP32-S3 beserta antarmuka pengguna pada perangkat, \textit{pipeline} preprocessing sinyal ECG pada mikrokontroler, implementasi model TinyML untuk deteksi aritmia secara \textit{on-device}, serta mekanisme bufer lokal menggunakan protokol MQTT untuk menjaga resiliensi data saat koneksi terputus.

\noindent
\textbf{\makebox[1.5cm][l]{BAB 4}EKSPERIMEN DAN ANALISIS}\\
\hspace*{1.5cm} Bab ini memaparkan parameter eksperimen, karakteristik data, tempat dan waktu ujicoba, spesifikasi peralatan, serta hasil eksperimen yang disusun per tujuan, yaitu jalur preprocessing \textit{real-time} di mikrokontroler, jalur model TinyML INT8 di ESP32-S3, dan jalur \textit{pipeline edge-to-cloud} yang resilien. Bab ini ditutup dengan analisis ketercapaian setiap tujuan, perbandingan dengan penelitian terkait, dan keterbatasan sistem.

\noindent
\textbf{\makebox[1.5cm][l]{BAB 5}PENUTUP}\\
\hspace*{1.5cm} Bab ini menyimpulkan hasil penelitian terhadap setiap tujuan beserta temuan teknis utamanya, kemudian memberikan saran pengembangan di luar lingkup penelitian ini.
